%------------------------------------------------------------------------------%
\section{Definição de Novas Funções}

Podemos definir novas funções como a soma de uma Série de Taylor, essa é uma
forma para definir as funções reais conhecidas para variáveis complexas,
como ilustrado na próxima subseção,
e também é como definimos várias funções especiais. Como exemplo vamos
definir uma das \emph{Funções de Bessel}\index{Função!de Bessel},
que são as soluções da equação de Bessel
\begin{equation}
  \label{eq:bessel}
  x^2 y'' + xy' + (x^2-\alpha^2)y = 0 \;.
\end{equation}

%------------------------------------------------------------------------------%
\begin{example}
  Exiba a Função de Bessel de Primeira Espécie de Ordem Zero, $J_0(x)$.
  \solution
  Essa função é uma solução da equação de Bessel~\eqref{eq:bessel} com
  $\alpha=0$
  \[
  \mathcal{L}[y] = x^2 y'' + xy' + x^2y = 0 \;.
  \]
  Vamos buscar a solução com a forma
  \[
  y(x) = \sum_{n=0}^\infty a_n x^n \;.
  \]
  Calculando as derivadas de $y$ obtemos
  \begin{align*}
    y'(x)  & = \sum_{n=1}^\infty a_n n x^{n-1} \\[3mm]
    y''(x) & = \sum_{n=2}^\infty a_n n(n-1) x^{n-2} \;.
  \end{align*}
  Substituindo na equação temos
  \begin{align*}
    \mathcal{L}[y] & = x^2 y'' + xy' + x^2y \\[3mm]
    & = x^2 \sum_{n=2}^\infty a_n n(n-1) x^{n-2}
    + x   \sum_{n=1}^\infty a_n n x^{n-1}
    + x^2 \sum_{n=0}^\infty a_n x^n
    \\[3mm]
    & = \sum_{n=2}^\infty a_n n (n-1) x^n
    + \sum_{n=1}^\infty a_n n       x^n
    + \sum_{n=0}^\infty a_n         x^{n+2} \;.
  \end{align*}
  Queremos agrupar todos os termos em um único somatório,
  para isso fazemos uma mudança de índice no último
  para igualar as potências de $x$
  \[
  \mathcal{L}[y] = \sum_{n=2}^\infty a_n n (n-1) x^n
  + \sum_{n=1}^\infty a_n n       x^n
  + \sum_{n=2}^\infty a_{n-2}     x^n \;.
  \]
  Precisamos agora fazer o índice do segundo somatório começar em $n=2$,
  conseguimos isso movendo o termo $n=1$ para fora
  \begin{align*}
    \mathcal{L}[y] & =        \sum_{n=2}^\infty a_n n (n-1) x^n
    + a_1x + \sum_{n=2}^\infty a_n n       x^n
    +        \sum_{n=2}^\infty a_{n-2}     x^n \\[3mm]
    & = a_1x + \sum_{n=2}^\infty \big( a_n n (n-1) + a_n n + a_{n-2} \big) x^n \\[3mm]
    & = a_1x + \sum_{n=2}^\infty \big( a_n n^2 + a_{n-2} \big) x^n \;.
  \end{align*}
  Para que
  \(
  \mathcal{L}[y] = 0
  \)
  precisamos que o coeficiente
  de cada potência de $x$ seja zero. Ao analisarmos o
  coeficiente de $x$ verificamos que
  \[
  a_1 = 0 \;.
  \]
  Para as demais potências, $x^n$, temos a condição
  \[
  a_n n^2 + a_{n-2} = 0
  \]
  que produz a relação de recorrência
  \[
  a_n = -\frac{a_{n-2}}{n^2} \;.
  \]
  Usando essa relação e a condição $a_1 = 0$ concluímos que todos
  os coeficientes ímpares são nulos
  \[
  a_n = 0 \qquad n = 1, 3, 5, \dots
  \]
  Para encontrarmos os coeficientes pares fazemos $n=2k$
  \[
  a_{2k} = -\frac{a_{2k-2}}{(2k)^2} \qquad k = 1, 2, 3, \dots
  \]
  Observando que
  \begin{align*}
    a_2 & = -\frac{a_0}{2^2} \\[3mm]
    a_4 & = -\frac{a_2}{4^2}
    = -\frac{1}{4^2}\left(-\frac{a_0}{2^2}\right)
    =  \frac{a_0}{2^6} \\[3mm]
    a_6 & = -\frac{a_4}{6^2}
    = -\frac{1}{6^2}\left(\frac{a_0}{2^6}\right)
    = -\frac{a_0}{2^6(3\cdot 2)^2} \\[3mm]
    a_8 & = -\frac{a_6}{8^2}
    = -\frac{1}{8^2}\left(-\frac{a_0}{2^6(3\cdot 2)^2}\right)
    =  \frac{a_0}{2^8(4\cdot 3\cdot 2)^2} \;.
  \end{align*}
  Generalizando temos
  \[
  a_{2k} = \frac{(-1)^k}{\,2^{2k}\,(k!)^2\,} a_0 \qquad k = 1, 2, 3, \dots
  \]
  Podemos agora escrever
  \[
  y(x) = a_0 \sum_{k=0}^\infty \frac{\,(-1)^k\,}{\,2^{2k}\,(k!)^2\,} \, x^{2k} \;.
  \]
  A \emph{Função de Bessel de Primeira Espécie de Ordem Zero} é definida como
  \[
  J_0(x) = \sum_{k=0}^\infty \frac{\,(-1)^k\,}{\,2^{2k}\,(k!)^2\,} \, x^{2k} \;.
  \]
  %----------------------------------------------------------------------------%
  Vamos agora analisar a convergência dessa série usando o Teste da Razão
  \[
  \abs{\frac{u_{k+1}}{u_k}}
  = \abs{ \frac{\,(-1)^{k+1} x^{2k+2}\,}{\,2^{2k+2}\,((k+1)!)^2\,}\;
    \frac{\,2^{2k}\,(k!)^2\,}{\,(-1)^k x^{2k}\,} }
  = \frac{x^2}{\,2^2 (k+1)^2\,}
  \]
  Quando $k\to\infty$ essa expressão vai para zero independentemente do valor
  de $x$, portanto a série converge para qualquer $x\in\R$.
  %----------------------------------------------------------------------------%
\end{example}
%------------------------------------------------------------------------------%

\begin{figure}
  \begin{center}%
    \vspace{-0.2\baselineskip}%
    \scalebox{0.9}{% Title: gl2ps_renderer figure
% Creator: GL2PS 1.4.2, (C) 1999-2020 C. Geuzaine
% For: Octave
% CreationDate: Tue Sep 28 14:13:47 2021
\setlength{\unitlength}{1pt}
\begin{picture}(0,0)
\includegraphics[width=17cm,height=7cm]{fig_taylor_bessel}
\end{picture}%
\begin{picture}(481,198)(0,0)
\fontsize{10}{0}\selectfont\put(88.9306,67.8458){\makebox(0,0)[t]{\textcolor[rgb]{0.15,0.15,0.15}{{5}}}}
\fontsize{10}{0}\selectfont\put(159.024,67.8458){\makebox(0,0)[t]{\textcolor[rgb]{0.15,0.15,0.15}{{10}}}}
\fontsize{10}{0}\selectfont\put(229.117,67.8458){\makebox(0,0)[t]{\textcolor[rgb]{0.15,0.15,0.15}{{15}}}}
\fontsize{10}{0}\selectfont\put(299.21,67.8458){\makebox(0,0)[t]{\textcolor[rgb]{0.15,0.15,0.15}{{20}}}}
\fontsize{10}{0}\selectfont\put(369.303,67.8458){\makebox(0,0)[t]{\textcolor[rgb]{0.15,0.15,0.15}{{25}}}}
\fontsize{10}{0}\selectfont\put(439.396,67.8458){\makebox(0,0)[t]{\textcolor[rgb]{0.15,0.15,0.15}{{30}}}}
\fontsize{10}{0}\selectfont\put(13.8417,20.4678){\makebox(0,0)[r]{\textcolor[rgb]{0.15,0.15,0.15}{{-0,5}}}}
\fontsize{10}{0}\selectfont\put(13.8417,130.185){\makebox(0,0)[r]{\textcolor[rgb]{0.15,0.15,0.15}{{0,5}}}}
\fontsize{10}{0}\selectfont\put(13.8417,185.044){\makebox(0,0)[r]{\textcolor[rgb]{0.15,0.15,0.15}{{1}}}}
\fontsize{10}{0}\selectfont\put(25.8468,190.53){\makebox(0,0)[l]{\textcolor[rgb]{0,0,0}{{$J_0(x)$}}}}
\fontsize{10}{0}\selectfont\put(166.033,163.101){\makebox(0,0)[l]{\textcolor[rgb]{0,0,0}{{$P_{20}$}}}}
\fontsize{10}{0}\selectfont\put(271.173,163.101){\makebox(0,0)[l]{\textcolor[rgb]{0,0,0}{{$P_{40}$}}}}
\fontsize{10}{0}\selectfont\put(376.312,163.101){\makebox(0,0)[l]{\textcolor[rgb]{0,0,0}{{$P_{60}$}}}}
\fontsize{10}{0}\selectfont\put(45.4729,31.4396){\makebox(0,0)[r]{\textcolor[rgb]{0,0,0}{{$P_{2}$}}}}
\fontsize{10}{0}\selectfont\put(104.351,31.4396){\makebox(0,0)[l]{\textcolor[rgb]{0,0,0}{{$P_{10}$}}}}
\fontsize{10}{0}\selectfont\put(215.098,31.4396){\makebox(0,0)[l]{\textcolor[rgb]{0,0,0}{{$P_{30}$}}}}
\fontsize{10}{0}\selectfont\put(320.238,31.4396){\makebox(0,0)[l]{\textcolor[rgb]{0,0,0}{{$P_{50}$}}}}
\fontsize{10}{0}\selectfont\put(425.377,31.4396){\makebox(0,0)[l]{\textcolor[rgb]{0,0,0}{{$P_{70}$}}}}
\end{picture}
}%
    \vspace{-\baselineskip}%
  \end{center}
  \caption{Comparação da função de Bessel com seus polinômios de Taylor.}
  \label{fig:funcao_tayplor_bessel}
\end{figure}

A Figura~\ref{fig:funcao_tayplor_bessel} a seguir
mostra o gráfico da função $J_0(x)$ e alguns dos seus Polinômios de Taylor.

%------------------------------------------------------------------------------%
